\chapter{Spin, Exchange, and the Two Minus Signs}

These notes follow Leonard Susskind's fifth lecture in the supplementary Advanced Quantum Mechanics sequence, with the lecture text polished into a coherent chapter and curated by LazyingArt LLC. The lecture begins with chemistry rather than formal many-body theory, because atomic shell filling is where spin and exclusion first appear together in an unmistakable way. From that empirical starting point Susskind moves to the wavefunctions of identical particles, derives bosonic and fermionic exchange symmetry, shows how the Pauli principle emerges from antisymmetry, and then turns to the second minus sign: the sign acquired by half-integer-spin states under a full \(2\pi\) rotation.

\section{Atomic Shells as the First Clue}

Susskind begins with a reminder. The fermionic character of the electron seems to show itself in two ways. First, the electron has spin \(1/2\), and therefore two spin states. Second, electrons obey the Pauli exclusion principle. Historically these facts appeared together in chemistry, so the lecture starts there.

In the simple noninteracting Coulomb-shell picture, the state of an electron is not specified by its spatial orbital alone. Spin counts as part of the state. This is why two electrons may occupy the lowest spatial orbital of helium: they do so only because their spin labels differ. In the rough shell-model language, one electron is spin up and the other is spin down. A third electron, as in lithium, cannot be placed into that same full state and must go into a higher shell.

The lecture keeps this shell model deliberately elementary. The point is not accurate atomic structure, but a counting argument. Orbital angular momentum \(l\) supplies \(2l+1\) spatial states:
\begin{equation}
l=0 \Rightarrow 2l+1=1,\qquad
l=1 \Rightarrow 2l+1=3,\qquad
l=2 \Rightarrow 2l+1=5.
\end{equation}
In spectroscopic notation these are the \(S\), \(P\), and \(D\) columns.

\begin{figure}
\begin{tikzpicture}
\draw (0,0)--(0,4.6);
\draw (2,0)--(2,4.6);
\draw (4,0)--(4,4.6);

\draw (-0.18,0.7)--(0.18,0.7);
\draw (-0.18,1.5)--(0.18,1.5);
\draw (-0.18,2.3)--(0.18,2.3);
\draw (-0.18,3.1)--(0.18,3.1);
\draw (-0.18,3.9)--(0.18,3.9);

\draw (1.82,2.3)--(2.18,2.3);
\draw (1.82,3.1)--(2.18,3.1);
\draw (1.82,3.9)--(2.18,3.9);

\draw (3.82,3.1)--(4.18,3.1);
\draw (3.82,3.9)--(4.18,3.9);

\node at (0,-0.35) {$S$};
\node at (2,-0.35) {$P$};
\node at (4,-0.35) {$D$};

\node at (0,-0.8) {$1$};
\node at (2,-0.8) {$3$};
\node at (4,-0.8) {$5$};
\end{tikzpicture}
\end{figure}

The lowest shell is an \(S\)-wave shell, so it has a single spatial orbital. The next shell contains one \(2s\) orbital and three \(2p\) orbitals, for a total of four spatial states:
\begin{equation}
N_{\mathrm{spatial}}(n=2)=1_{2s}+3_{2p}=4.
\end{equation}
Since the electron also carries two spin states, the number of electron states in this shell doubles:
\begin{equation}
N_{\mathrm{electron}}(n=2)=2\,N_{\mathrm{spatial}}(n=2)=8.
\end{equation}
That is the occupancy count Susskind uses for the first excited shell. The shell picture then explains, at least heuristically, why neon closes this shell and why sodium again has one outer electron.

At this point Susskind interrupts the chemistry himself. He does not want to do chemistry for its own sake; he wants to remind us of a striking empirical pairing. Half-integer spin and exclusion seem to go together. But must they?

To sharpen the question he contrasts electrons with photons. Photons certainly do not obey exclusion. Classical electromagnetic waves are just situations in which very many photons occupy the same state. Thus integer-spin particles behave in the opposite way from electrons, and the lecture turns from chemistry to the mathematics of identical particles.

\section{Many-Particle Wavefunctions and Identical Particles}

For a system of \(n\) particles, a convenient basis is obtained by specifying all of their positions:
\begin{equation}
|x_1,x_2,\ldots,x_n\rangle.
\end{equation}
Here \(x_i\) labels particle \(i\). It does not mean the \(x\)-component of position; the three spatial coordinates are being suppressed into one symbol.

If the abstract state of the system is denoted by \(|\Psi\rangle\), its coordinate-space wavefunction is
\begin{equation}
\psi(x_1,x_2,\ldots,x_n)=\langle x_1,x_2,\ldots,x_n\mid \Psi\rangle.
\end{equation}
Its squared magnitude is the joint probability density:
\begin{equation}
|\psi(x_1,x_2,\ldots,x_n)|^2.
\end{equation}
This is the probability density for finding particle \(1\) at \(x_1\), particle \(2\) at \(x_2\), and so on.

The new issue arises when the particles are identical. Is the state with particle \(1\) at \(x_1\) and particle \(2\) at \(x_2\) the same physical state as the state with particle \(1\) at \(x_2\) and particle \(2\) at \(x_1\)? For distinguishable particles that question would be trivial, but for electrons it is not. Electrons do not come in observable colors; one cannot physically tag them as ``the blue electron'' and ``the red electron.'' The lecture therefore asks for the correct relation between the two labelings.

A further ingredient is that quantum mechanics is insensitive to an overall phase. If two state vectors differ only by a common phase factor, they describe the same physics. That fact opens the door to a slightly weaker statement than literal equality.

\section{Exchange, Bosons, and Fermions}

Suppose that interchanging two identical particles changes the wavefunction only by a phase:
\begin{equation}
\psi(x_1,x_2)=e^{i\phi}\,\psi(x_2,x_1).
\end{equation}
Now interchange them again. After two exchanges the particles are back where they started, so the state must return to itself:
\begin{align}
\psi(x_1,x_2)
&=e^{i\phi}\psi(x_2,x_1) \\
&=e^{2i\phi}\psi(x_1,x_2).
\end{align}
Therefore
\begin{equation}
e^{2i\phi}=1.
\end{equation}
In the three-dimensional setting of this lecture, Susskind keeps the two resulting possibilities:
\begin{equation}
e^{i\phi}=+1
\qquad \text{or} \qquad
e^{i\phi}=-1.
\end{equation}

\begin{definition}
A bosonic wavefunction is symmetric under exchange:
\begin{equation}
\psi_B(x_1,x_2)=+\psi_B(x_2,x_1).
\end{equation}
A fermionic wavefunction is antisymmetric under exchange:
\begin{equation}
\psi_F(x_1,x_2)=-\psi_F(x_2,x_1).
\end{equation}
\end{definition}

\begin{remark}
As Susskind emphasizes in passing, this is the three-dimensional story. In one and two spatial dimensions the exchange structure is more subtle, so the reduction to just \(\pm 1\) should not be stated as a universal theorem of all dimensions.
\end{remark}

The important conceptual point is that the Pauli principle is no longer an extra sentence appended to quantum mechanics. It is beginning to look like a statement about the symmetry type of the many-particle wavefunction.

\section{Two-Particle States and the Pauli Principle}

To see what the exchange rule means concretely, Susskind turns to two one-particle orbitals. On the board he labels them \(\psi_0,\psi_1,\psi_2\), with \(\psi_0\) the lowest state. If particle \(1\) is placed in \(\psi_0\), its factor is \(\psi_0(x_1)\).

\paragraph{Two particles in the same state.}

Start with the product
\begin{equation}
\psi_{00}(x_1,x_2)=\psi_0(x_1)\psi_0(x_2).
\end{equation}
Because wavefunctions are ordinary numerical functions and not operators, the order of multiplication does not matter. Thus
\begin{equation}
\psi_{00}(x_2,x_1)=\psi_0(x_2)\psi_0(x_1)=\psi_0(x_1)\psi_0(x_2)=\psi_{00}(x_1,x_2).
\end{equation}
So this state is symmetric. It is therefore acceptable for bosons, but not for fermions.

\paragraph{Two particles in different states.}

Now place one particle in \(\psi_0\) and the other in another one-particle state \(\psi_1\):
\begin{equation}
\psi^{\mathrm{prod}}_{01}(x_1,x_2)=\psi_0(x_1)\psi_1(x_2).
\end{equation}
After exchanging the particles,
\begin{equation}
\psi^{\mathrm{prod}}_{01}(x_2,x_1)=\psi_0(x_2)\psi_1(x_1).
\end{equation}
In general the product state is neither equal to its exchanged version nor to its negative. By itself it is therefore not an allowed bosonic or fermionic wavefunction.

The cure is to build linear combinations. For two distinct one-particle states the symmetrized and antisymmetrized combinations are
\begin{align}
\psi_B^{(01)}(x_1,x_2)
&=\frac{1}{\sqrt2}\Big[\psi_0(x_1)\psi_1(x_2)+\psi_0(x_2)\psi_1(x_1)\Big], \\
\psi_F^{(01)}(x_1,x_2)
&=\frac{1}{\sqrt2}\Big[\psi_0(x_1)\psi_1(x_2)-\psi_0(x_2)\psi_1(x_1)\Big].
\end{align}
The first is unchanged by exchange; the second changes sign.

Now comes the lecture's key payoff. If the two one-particle states are actually the same, the fermionic combination vanishes:
\begin{align}
\psi_F^{(00)}(x_1,x_2)
&\propto \psi_0(x_1)\psi_0(x_2)-\psi_0(x_2)\psi_0(x_1) \\
&=0.
\end{align}
This is Pauli exclusion in explicit mathematical form. The exclusion principle is not added by hand; it is the statement that an antisymmetric two-particle wavefunction vanishes when the two full one-particle states coincide.

The bosonic story is the opposite. When both particles occupy the same orbital, the bosonic state survives. After normalization it is simply the product \(\psi_0(x_1)\psi_0(x_2)\).

\paragraph{Superposition versus product.}

A useful clarification appears in the student discussion. Adding one-particle wavefunctions describes one particle in a superposition of states:
\begin{equation}
\psi_{\mathrm{one\ particle}}=\alpha\,\psi_0+\beta\,\psi_1.
\end{equation}
By contrast, a two-particle product state is built by multiplication:
\begin{equation}
\psi_{\mathrm{two\ particles}}(x_1,x_2)=\psi_0(x_1)\psi_1(x_2).
\end{equation}
Only after one has built the two-particle product does one symmetrize or antisymmetrize it.

This is also where Susskind agrees with a student that the signs \(\pm 1\) form a representation of the permutation group. In a cleaned-up shorthand one may write
\begin{equation}
P_{12}\psi_{\pm}=\pm \psi_{\pm},
\end{equation}
but the actual content is still the explicit transformation of the arguments.

\section{Including Spin in the Particle Label}

Until this point the lecture has deliberately ignored spin. Now Susskind puts it back in. The easiest way is not to change the formalism, but to enlarge the particle label. A one-particle wavefunction becomes
\begin{equation}
\psi(x,\sigma),
\end{equation}
where \(\sigma\) is a discrete spin label. In the lecture this is tied to a chosen \(\sigma_z\) basis, so one may think of \(\sigma=\pm 1\) as the two possible spin labels in that basis. The probability density is
\begin{equation}
|\psi(x,\sigma)|^2,
\end{equation}
interpreted as the probability density for finding the particle at position \(x\) with the corresponding spin value.

For many particles one now writes the wavefunction as a function of all the full labels,
\begin{equation}
\psi(x_1,\sigma_1;\,x_2,\sigma_2;\,\ldots).
\end{equation}
Exchange must act on everything belonging to particle \(1\) and particle \(2\), not just on their positions:
\begin{align}
\psi_B(x_1,\sigma_1;\,x_2,\sigma_2)
&=+\psi_B(x_2,\sigma_2;\,x_1,\sigma_1), \\
\psi_F(x_1,\sigma_1;\,x_2,\sigma_2)
&=-\psi_F(x_2,\sigma_2;\,x_1,\sigma_1).
\end{align}
Thus the Pauli principle must be read correctly: two fermions cannot occupy the same full state, including spin.

In particular,
\begin{equation}
\psi_F(x,\sigma;\,x,\sigma)=0.
\end{equation}
But two electrons may still share the same spatial orbital if their spin labels differ, because \((x,\uparrow)\) and \((x,\downarrow)\) are not the same one-particle state. This is the correct mathematical statement behind the helium example with which the lecture began.

At this point Susskind also endorses a useful interpretation: the vanishing of the antisymmetric same-state combination may be viewed as a kind of destructive interference between the two exchanged configurations. He is careful, however, not to let that interpretation replace the symmetry rule itself.

\section{Rotation by \(2\pi\) and the Second Minus Sign}

Having finished with exchange, Susskind turns to the second minus sign. Consider a one-particle state, and imagine rotating it about the \(z\)-axis by an angle \(\theta\). The total angular momentum, including both orbital and spin parts, is denoted by \(J\), and its \(z\)-component is the generator of these rotations.

In the lecture the generator statement is expressed in differential form:
\begin{equation}
J_z\,\psi(\theta)=-i\,\frac{\partial}{\partial\theta}\psi(\theta).
\end{equation}
If the state has definite \(J_z\)-eigenvalue \(m\), then
\begin{equation}
J_z\,\psi=m\,\psi.
\end{equation}
Combining the two equations gives a first-order differential equation for the angle-dependent state:
\begin{equation}
-i\,\frac{\partial}{\partial\theta}\psi(\theta)=m\,\psi(\theta).
\end{equation}
Its solution is
\begin{equation}
\psi(\theta)=e^{im\theta}\psi(0).
\end{equation}
Therefore a full rotation produces
\begin{equation}
\psi(2\pi)=e^{i2\pi m}\psi(0).
\end{equation}

Now use the angular momentum spectrum already known from the commutation-relation analysis. The allowed values of \(m\) are either integers or half-integers. Thus
\begin{equation}
m\in \mathbb{Z} \quad \Rightarrow \quad e^{i2\pi m}=+1,
\end{equation}
while
\begin{equation}
m\in \mathbb{Z}+\frac12 \quad \Rightarrow \quad e^{i2\pi m}=-1.
\end{equation}
So integer-spin states come back to themselves after a \(2\pi\) rotation, whereas half-integer-spin states come back only up to a minus sign. After a \(4\pi\) rotation all of them return:
\begin{equation}
\psi(4\pi)=\psi(0).
\end{equation}

\begin{remark}
In more standard operator notation one often writes \(R_z(\theta)=e^{-i\theta J_z}\). The lecture instead emphasizes the angle-differentiation rule \(J_z=-i\,\partial/\partial\theta\) acting on the rotated state. The sign convention in the exponent is therefore best treated cautiously. The invariant content is the \(2\pi\) behavior: integer spin gives \(+1\), half-integer spin gives \(-1\), and \(4\pi\) gives the identity.
\end{remark}

This is the second minus sign in the lecture: not the sign from exchanging two fermions, but the sign acquired by rotating a half-integer-spin state by \(2\pi\).

\section{Topology, Spin-Statistics, and the Observability of the Sign}

Susskind next changes style. The algebra is not abandoned, but it is supplemented by topology and physical pictures.

\paragraph{The belt and the topology of \(4\pi\).}

The cup trick and the belt trick are used to display a basic topological fact about the rotation group. A path in the space of rotations corresponding to a \(2\pi\) rotation is not trivial, but the corresponding \(4\pi\) path is. In the language of the lecture, a \(2\pi\) twist produces a tangle, while a second \(2\pi\) twist allows the tangle to be unwound.

This is not yet the spin-statistics theorem. But it is meant to show that there really is mathematics behind the peculiar statement that a \(2\pi\) rotation is not quite the identity for fermions, while a \(4\pi\) rotation is.

\paragraph{Exchange plus rotation.}

The next step is heuristic and should remain labeled as such. Susskind uses a belt and then a worldline-style picture to suggest that an exchange of two particles and a \(2\pi\) rotation of one particle are topologically linked. In the worldline version one creates two particle-antiparticle pairs, rotates one particle by \(2\pi\), exchanges another pair, and then deforms the whole history into a trivial creation-annihilation process. If that deformation is topologically legal, then the phase from exchange must cancel the phase from the \(2\pi\) rotation.

The symbolic moral is that
\begin{equation}
\text{exchange phase} \times \text{rotation phase} = 1
\end{equation}
for the combined topologically trivial process. This is not Pauli's proof, and Susskind does not present it as a rigorous theorem. He presents it as a suggestive explanation of why half-integer spin should go together with fermionic exchange, and integer spin with bosonic exchange.

He also stresses that this heuristic is broader than the original relativistic proof. It is meant to apply not only to elementary particles in relativistic quantum field theory, but more generally to objects that can be created and annihilated: solitons, lumps of fields, monopoles, and similar excitations.

\paragraph{Why a global sign is usually invisible.}

The lecture closes by asking whether the minus sign from a \(2\pi\) rotation can actually be observed. If one traps an electron in a magnetic field, rotates it adiabatically by \(2\pi\), and then compares that experiment to one in which no such rotation was performed, the result is null. The sign is only global. Since probabilities involve \(\psi\psi^\ast\), the overall minus sign disappears.

\paragraph{Relative phase and interference.}

The sign becomes observable only if it is turned into a relative phase between coherent branches. Susskind therefore imagines a one-electron interferometer. One electron is sent into a beam splitter, so that the state becomes a superposition of two separated wave packets:
\begin{equation}
\psi(x)=\psi_{\mathrm{L}}(x)+\psi_{\mathrm{R}}(x).
\end{equation}
The electron is still a single electron. The wavefunction has simply been split into two coherent branches.

Now let only one branch execute a \(2\pi\) spin rotation while the other branch is left alone. The state changes to
\begin{equation}
\psi(x)\longrightarrow \psi_{\mathrm{L}}(x)-\psi_{\mathrm{R}}(x).
\end{equation}
When the branches are recombined, the interference pattern is different. The plus sign and the minus sign lead to different constructive and destructive interference. In particular, at a symmetry point one may have a node in the minus case and a nonzero amplitude in the plus case. Thus the sign of the \(2\pi\) rotation becomes experimentally meaningful because it is no longer global.

Susskind then describes a more realistic version of the experiment. Instead of literally turning a box slowly by hand, one arranges a magnetic field in only one branch so that the spin there precesses through one full \(2\pi\) rotation while the other branch does not. The actual experimental details are not the main point. The essential point is this: a global sign is invisible, but a relative sign between coherent branches is observable.

\section{Summary}

The lecture begins with atomic shell filling because that is where spin and exclusion first seem to appear together. It then recasts exclusion as a statement about the symmetry of many-particle wavefunctions: bosons are symmetric, fermions are antisymmetric, and the Pauli principle follows from the vanishing of the antisymmetric same-state combination. After spin is restored as part of the particle label, the same symmetry rule says that fermions cannot occupy the same full state, including spin. Only then does Susskind turn to the second minus sign, deriving the \(-1\) acquired by half-integer-spin states under a \(2\pi\) rotation from the generator \(J_z\). The closing heuristic pictures and interferometer thought experiments are there to make one final point: the minus sign from exchange and the minus sign from \(2\pi\) rotation are not unrelated curiosities, and under the right interference conditions the latter is not merely formal but observable.